\section{The parity presentation and the component-chain ideal for the changed tuple}

This section uses the actual tuple, rings and completed local functor of
AC-G94--AC-G101.  In particular, the residue characteristic is
\(p=269\), the residual characters are
\((\chi^j\gamma,1,\chi^j\gamma^{-1})\), and the normalized
matrix at the distinguished element is
\[
 r(g_0)-r(g_0^\star)=\operatorname{diag}(a,0,-a),
 \qquad a\in B^\times .
 \tag{G105.1}
\]
No replacement of this tuple by powers of a single character is made.
Let \(x\) be a closed point over \(\mathfrak m_B\), and write
\(C=C_x\), with coefficient DVR \(\mathcal O'\), for its
completion from AC-G101.  We use indices \(0,1,2\).  Let
\(\mathfrak e\) be the ideal generated by all entries joining
the middle index to either endpoint.  These are exactly the
entries of odd weight for \(w=(1,0,-1)\).  The notation
\(\mathfrak e_A\) denotes the corresponding ideal in \(A\).

\begin{theorem}[AC-G105-parity: the actual even--odd presentation]
\label{AC-G105-parity}
If \(\mathfrak e_A\subset x\), there are nonnegative integers
\(a_+,b_+,g\), a continuous surjection
\[
 \begin{split}
 q:S=\mathcal O'[[X_1,\ldots,X_{a_+},Y_1,\ldots,Y_g]]&\longrightarrow C,\qquad
 a_+=b_++3d+1,\\
 \ker q&=(F_1,\ldots,F_{b_+},G_1,\ldots,G_g),
 \end{split}
 \tag{G105.2}
\]
such that each \(F_i\) is even and each \(G_i\) is odd under
\(X\mapsto X,\ Y\mapsto-Y\).  Every relation has its linear
coefficients in \(\mathfrak m_{\mathcal O'}\).  All entries of
the universal representation have lifts of their prescribed
weight parity, and every element of \(B\) has an even lift.
Moreover
\[
                         \mathfrak e=(q(Y_1),\ldots,q(Y_g)).
 \tag{G105.3}
\]
Thus this is the literal numerical and parity content of
\texttt{Space.Loc.EPres} for the actual tuple.
\end{theorem}

\begin{proof}
\emph{1. The actual involution and the local flag.}
Put \(P=\operatorname{diag}(-1,1,-1)\).  The weight grading
gives an \(\mathcal O\)-involution \(\iota\) of \(A\),
fixing \(B\), with
\[
                       \iota(r(s))=P r(s)P.
 \tag{G105.4}
\]
It fixes \(x\) and acts identically on its residue field:
the odd generators vanish there and the even generators are
fixed.  It therefore extends to an \(\mathcal O'\)-involution
of \(C\).  The residual representation \(\sigma\) commutes
with \(P\).

The local ordered flag can be chosen compatibly with this
involution, without assuming an unproved equivariant flag
lifting theorem.  Choose the separating element \(s_v\) used
in AC-G101.  Its three labeled eigenvalues
\(\kappa_{v,i}(s_v)\) lie in \(B\), and their differences
are units.  Its Hensel projectors are
\[
 E_{v,i}=\prod_{h\ne i}
       \frac{r(s_v)-\kappa_{v,h}(s_v)}
            {\kappa_{v,i}(s_v)-\kappa_{v,h}(s_v)}.
 \tag{G105.5}
\]
They satisfy \(\iota(E_{v,i})=PE_{v,i}P\).  Since the
residual matrices at \(x\) are upper or lower triangular
with the labeled diagonal, the \(i\)-th coordinate of
\(E_{v,i}e_i\) is a unit.  Divide each column by that
coordinate, obtaining an invertible matrix \(V_v\).  Its
reduction is triangular with diagonal entries one, and
\[
                         \iota(V_v)=PV_vP.
 \tag{G105.6}
\]
These columns diagonalize \(r(s_v)\).  The existing ordered
flag is the succession of the labeled eigenspaces, by the
uniqueness in AC-G101.  Thus \(V_v^{-1}r(s)V_v\) is upper
triangular for every local \(s\).  This also proves that the
local Borel Lie algebras respect the parity decomposition.

\emph{2. The two actual Euler characteristics.}
Let \(W=M_3(k')\) be the actual adjoint representation of
\(\Gamma\) in AC-G101, and split it by conjugation by \(P\):
\(W=W^+\oplus W^-\).  The first summand consists of the
endpoint block and the middle diagonal entry; the second
consists of the four cross entries.  Write
\(H^i_\pm\) for the continuous cone with global coefficient
\(W^\pm\) and local quotients
\(W^\pm/(\mathfrak b_v\cap W^\pm)\), and put
\(h_i^\pm=\dim_{k'}H^i_\pm\).  The dimensions are
\[
\begin{array}{c|cc}
 &+&-\\ \hline
 W&5&4\\
 \mathfrak b_v\cap W^\pm&4&2\\
 (W^\pm)^{c_w=1}\quad\text{at a real place}&1&2.
\end{array}
 \tag{G105.7}
\]
For the last row, the matrix defining the real involution is
invertible symmetric and preserves the endpoint--middle
decomposition.  Multiplication by that matrix identifies
the invariant equation with a skew-symmetric matrix equation.
The endpoint skew block has dimension one, the middle skew
block has dimension zero, and the two cross positions give
dimension two.  This argument works at every real conjugate,
since \(\sigma(G)\) preserves the same decomposition.

The actual dual vanishing (G101.18), with \(c_p\) denoting
the action on \(\mu_p\), applies to each direct summand.
The EC identity of AC-G96, including its absence of a
degree-three term, therefore gives
\[
 \begin{split}
 h_1^+&=h_2^++h_0^++3d,\\
 h_1^-&=h_2^-+h_0^-.
 \end{split}
 \tag{G105.8}
\]
Here the dual exclusion is the actual local calculation for
both cyclotomic conventions in AC-G101; it is not the old
one-character hypothesis imported unchanged.

Let \(T=T^+\oplus T^-\) be the relative tangent space of
\(C\).  The map \(T\to H^1_+\oplus H^1_-\) proved
surjective in AC-G101 commutes with parity.  Its restrictions
are surjective, by averaging with \((1\pm\iota)/2\).
The remaining normalization direction
\(X_0=\operatorname{diag}(1,0,-1)\) is even.  If it is
globally invariant, it supplies \(h_0^+\ge1\); if it is
not invariant, conjugation by
\(\operatorname{diag}(1+t,1,(1+t)^{-1})\), \(t^2=0\),
supplies a nonzero even tangent in the kernel.  Consequently
\[
 \dim T^+\ge h_2^++3d+1,
 \qquad \dim T^-\ge h_2^-.
 \tag{G105.9}
\]
All statements about flags and coboundaries here use the
vanishing \(H^0(G_{F_v},W/\mathfrak b_v)=0\) already
proved in (G101.17), so no extra local flag parameters have
been discarded.

\emph{3. The obstruction is equivariant, including complete targets.}
Choose an eigenbasis of the relative cotangent space and lift
it by averaging to even and odd elements of \(C\).  This
gives a minimal equivariant surjection
\[
 S_0=\mathcal O'[[X_1,\ldots,X_{a_+},Y_1,\ldots,Y_g]]
       \longrightarrow C,
 \quad a_+=\dim T^+,\quad g=\dim T^-.
 \tag{G105.10}
\]
Its kernel \(\mathfrak a\) is stable under parity.
An eigenbasis of \(\mathfrak a/\mathfrak m_{S_0}\mathfrak a\),
again lifted by averaging, gives even and odd minimal relations.

We verify the equivariance of the complete obstruction used
to bound these two numbers.  For an equivariant small extension
\(E\to D\) with kernel killed by \(\mathfrak m_E\), take
the continuous set lifts constructed in AC-G101.  Replace
each entry of prescribed parity \(\nu\) by
\((z+(-1)^\nu\iota_E(z))/2\).  Its image in \(D\) is
unchanged.  This operation is continuous, preserves zero
entries of upper-triangular matrices, and preserves
invertibility since the reduction is unchanged.  Use (G105.6)
for the local flag matrices.  Products and inverses of these
graded matrices have the required parities, including the
polarized semidirect multiplication.  Thus the multiplication
defects and local comparison defects defining (G101.22a)
are equivariant.  Their class in
\((H^2_+\oplus H^2_-)\otimes\ker(E\to D)\) is invariant.
The strict normalizations and the criterion that vanishing
of this class permits a lift are exactly those proved in
AC-G101; no new obstruction space is asserted.

Apply this to
\(S_0/\mathfrak m_{S_0}\mathfrak a\to C\).
The injection of the dual relation space into \(H^2\)
proved in AC-G101 is now an equivariant injection.  If
\(b_0\) and \(g_0\) count the even and odd minimal relations,
respectively, then
\[
            b_0\le h_2^+\le a_+-(3d+1),
            \qquad g_0\le h_2^-\le g.
 \tag{G105.11}
\]
Append zero even relations to reach \(b_+=a_+-(3d+1)\),
and zero odd relations to reach \(g\).  This gives (G105.2).
Minimality of (G105.10) says that the linear coefficients of
every relation lie in \(\mathfrak m_{\mathcal O'}\).
Equivariant averaging of a preimage under \(q\) gives the
claimed lifts of entries and of elements of \(B\).

Finally \(A/\mathfrak e_A\) is fixed pointwise by parity:
it is generated by \(B\) and the remaining even entries.
The same holds after localization, coefficient extension
and completion, since ideals in the complete Noetherian
local ring are closed.  Hence every odd element of \(C\)
belongs to \(\mathfrak e\).  Conversely every generating
entry of \(\mathfrak e\) is odd.  Every odd power series
belongs to the ideal generated by the \(Y_i\), so these
observations give (G105.3).
\end{proof}

\paragraph{The exact connectedness input.}
For a complete Noetherian local ring whose reduction is a
domain of dimension \(n\), cutting by \(r\) elements of
the maximal ideal leaves components of dimension at least
\(n-r\), connected to one another through intersections
of dimension at least \(n-r-1\).
The component statement is the height theorem and the
dimension formula for complete local rings.  For the chain
statement we use Grothendieck connectedness in the precise
form of Stacks Project, Lemma 52.14.1, tag
\texttt{0ECR}, \url{https://stacks.math.columbia.edu/tag/0ECR}.
Its entire printed proof was read.  Since an ideal generated
by \(r\) elements has cohomological dimension at most \(r\),
that lemma gives the assertion: a partition of the finite
set of components with all cross intersections too small
would disconnect their complement.  The cited duality and
formal-function inputs of that lemma remain published inputs.

\begin{lemma}[A determinant section and its component chains]
\label{G105-detsection}
Let \(\mathcal O'\) be a complete DVR and
\(S=\mathcal O'[[X_1,\ldots,X_a,Y_1,\ldots,Y_g]]\),
where \(g\ge1\).  Let \(F_1,\ldots,F_b\) belong to
\(\mathfrak m_S\), and let \(M\) be a \(g\)-by-\(g\)
matrix whose entries belong to \(\mathfrak m_S\).
If \(a=b+D\), the ideal
\[
                         (F_1,\ldots,F_b,MY,\det M)
 \tag{G105.12}
\]
has every minimal component of dimension at least \(D+1\);
any two such components are joined by a finite chain whose
successive intersections have dimension at least \(D\).
\end{lemma}

\begin{proof}
We give the algebra behind the exact source
\texttt{Detsec.detsec}, rather than treating the special
determinant cut as one more unconstrained equation.
For any Noetherian domain \(R\), introduce an independent
generic \(g\)-by-\(g\) matrix \(U\) and vector \(V\).
The radical of
\(J=(UV,\det U)\subset R[U,V]\) is prime.
Indeed put \(E=\operatorname{diag}(0,1,\ldots,1)\) and
consider the polynomial parametrization
\[
                  U=N E\operatorname{adj}(Z),
                  \qquad V=tZe_1.
 \tag{G105.13}
\]
Its kernel is prime, because its target is a polynomial
ring over a domain.  It is exactly \(\sqrt J\): over any
field, a pair satisfying \(UV=0,\det U=0\) has the
form (G105.13).  If \(V\ne0\), complete \(V\) to the
first column of an invertible \(Z\); if \(V=0\), instead
choose any nonzero kernel vector of \(U\).  Taking
\(N=(\det Z)^{-1}UZ\) and the appropriate scalar \(t\)
then gives the formula.  Apply this to the fraction field
of every prime quotient containing \(J\).  This proves
the asserted radical identity without an algebraically
closed-field or characteristic assumption.

The prime \(\sqrt J\) is homogeneous and does not contain
\(V_1\).  Homogeneous degree components show that its
extension to \(R[[U,V]]\) is also prime: a series belongs
to that extension exactly when all its homogeneous
components do.  To see this equivalence, choose finitely
many homogeneous generators and group the homogeneous
coefficient representations into formal series.  The
first nonzero homogeneous terms of a product prove
primality.  Thus the radical of the formal-series
incidence ideal is prime as well.  Its height is at most
\(g\).  A minimal prime of the \(g\) equations \(UV\)
contained in it already contains \(\det U\), since
\(\det U\,V_1\) lies in \((UV)\) and \(V_1\) is absent
from that prime.  The height theorem applies to these
\(g\) equations.

Use \(R=S\).  The irreducible reduction of
\(S[[U,V]]/(UV,\det U)\) consequently has dimension at
least \(1+a+g^2+g\).  Impose the \(g^2+g+b\) equations
\(U=M,\ V=Y,\ F=0\), all in the maximal ideal.
The connectedness input above gives component dimension
at least \(1+a-b=D+1\) and chain intersections of
dimension at least \(a-b=D\).  Formal substitution
identifies this quotient with (G105.12).
This proves the lemma, including nonreduced quotients.
\end{proof}

\begin{theorem}[AC-G105: the actual Sharp ideal and its precise dichotomy]
\label{AC-G105}
Assume that \(T_+\not\subset x\), and that the actual
triangular cut of AC-G98 has dimension less than \(3d\)
in \(C\).  There is an ideal \(\mathfrak s\subset C\)
with the following properties, which are the literal
algebraic content of \texttt{Space.Loc.Sharp}:
\begin{enumerate}
\item If \(\mathfrak e_A\not\subset x\), then
\(\mathfrak s=0\).
\item Every prime \(P\) not containing \(\mathfrak e\)
contains \(\mathfrak s\).
\item Every minimal prime of \(\mathfrak s\) has
dimension at least \(3d+2\).  Any two are joined by a
chain of minimal primes whose successive intersections
have dimension at least \(3d+1\).
\item At every such minimal prime, the representation
is off both triangular loci.  Either some odd entry
survives, or there is a continuous nonsplit extension
of its middle character by its endpoint block, with
the same even entries and both of the source ordered
local product identities.
\end{enumerate}
The nonsplitting holds after every injective map of the
prime quotient into a field.  Properness of
\(\mathfrak s\) is automatic if a prime not containing
\(\mathfrak e\) exists.  No unconditional existence of
such a prime is asserted here.
\end{theorem}

\begin{proof}
Set \(D=3d+1\).  If \(\mathfrak e_A\not\subset x\),
then \(\mathfrak e=C\).  Use the presentation of
AC-G101 with variable--relation excess at least \(D\).
The connectedness input gives the asserted chains for
\(\mathfrak s=0\).  No prime contains \(\mathfrak e\),
so only the first alternative of the dichotomy occurs.

Suppose \(\mathfrak e_A\subset x\), and take (G105.2).
If \(g=0\), (G105.3) says \(\mathfrak e=0\); set
\(\mathfrak s=C\).  The assertions about primes are
then vacuous, and the first property has a false
antecedent.  This edge case is intentional.  It does
not borrow the source's separate \texttt{gpos}
theorem, which invokes a nonzero-cocycle and
preconnectedness argument.

For \(g>0\), write the odd relations as
\[
             G_i=\sum_{j=1}^g M_{ij}Y_j,
             \qquad M_{ij}\in\mathfrak m_S,
             \qquad \mathfrak s=(q(\det M)).
 \tag{G105.14}
\]
One obtains this expression by assigning each monomial
of an odd series to its first odd variable.  The
constant coefficient of \(M_{ij}\) is a linear
coefficient of \(G_i\), hence lies in
\(\mathfrak m_{\mathcal O'}\).  The determinant-section
lemma proves the asserted chain and dimension bounds.
The adjugate identity gives
\(q(\det M)q(Y_j)=0\).  Thus every prime missing some
\(q(Y_j)\) contains \(\mathfrak s\), proving the second
property and the properness assertion.

The ideal \(T_+C\) is the unit ideal, since an element
of \(T_+\) is absent from \(x\).  The assumed bound
on \(C/T_-C\) is smaller than the dimension of every
minimal component of \(\mathfrak s\).  No such
component is in either triangular locus.
It remains to prove the nonsplit alternative at a
minimal prime \(P\) containing \(\mathfrak e\).
Put \(Q=C/P\), and denote the quotient map by
\(\pi:C\to Q\).

The matrix \(\pi q(M)\) has determinant zero.  Over
the domain \(Q\) there is a nonzero vector
\(\nu\in Q^g\) in its kernel: choose one over its
fraction field and clear a single denominator.
Define on \(S\)
\[
       \partial_\nu f
          =\sum_j \pi q\left(\frac{\partial f}{\partial Y_j}\right)
                    \nu_j.
 \tag{G105.15}
\]
It is a derivation relative to \(\pi q\).  It kills
every even series, since its derivative is odd and
all \(\pi q(Y_j)\) vanish.  It kills the odd
relations \(MY\), since the surviving term is
\(\pi q(M)\nu=0\).  The Leibniz rule therefore
makes it descend to a derivation
\(\partial:C\to Q\), zero on \(B\) and on all
even entries, with
\(\partial q(Y_j)=\nu_j\) not all zero.

For \(s\in G\) define
\[
 r_\partial(s)=
 \begin{pmatrix}
 \pi r_{00}(s)&\partial r_{01}(s)&\pi r_{02}(s)\\
 0&\pi r_{11}(s)&0\\
 \pi r_{20}(s)&\partial r_{21}(s)&\pi r_{22}(s)
 \end{pmatrix}.
 \tag{G105.16}
\]
Multiplicativity follows entry by entry from the
Leibniz rule, the vanishing of the odd entries in
\(Q\), and the vanishing of \(\partial\) on the
even entries.  The identity and inverse follow by
applying the same calculation to \(1\) and
\(ss^{-1}\).  Its middle quotient is the character
\(\pi r_{11}\), and its endpoint submodule is the
unchanged two-dimensional block.  The representation
is continuous in the quotient-adic topology:
\(\partial(\mathfrak m_C^{n+1})\subset
(\mathfrak m_CQ)^n\), by the Leibniz rule, and the
original entries are continuous.

For completeness, this construction preserves the
ordered local product identities on both the chosen
local subgroup and its star-conjugate.  On a matrix
whose odd entries vanish under \(\pi\) and whose
even entries are killed by \(\partial\), the
operation in (G105.16) commutes with multiplication.
It also commutes with subtracting a scalar killed
by \(\partial\).  All local characters belong to
\(B\), so this applies to each factor of the
ordered product.  On the star-conjugate subgroup
the characters occur in reversed order, exactly
as follows by transposing the original polarized
identity.  No local semisimplicity of
\(r_\partial\) is used.

Finally suppose that after an injective map
\(f:Q\to L\) to a field this extension split.
There would be a vector \((v_0,1,v_2)\) on which
every \(r_\partial(s)\) acts by its middle
character.  Compare that equation at \(g_0\)
and \(g_0^\star\).  Equation (G105.1) gives
\(f(\pi a)v_0=0\) and \(f(\pi a)v_2=0\):
the off-diagonal entries at those two elements
are equal.  Since \(a\) is a unit, \(v_0=v_2=0\).
The same equation for general \(s\) now gives
\(\partial r_{01}(s)=\partial r_{21}(s)=0\).
Polarization and substitution of \(s^\star\)
give the other two odd derivatives as zero.
The derivation then vanishes on their generated
ideal \(\mathfrak e\), by the Leibniz rule and
\(\pi\mathfrak e=0\).  This contradicts
\(\partial q(Y_j)=\nu_j\ne0\) for some \(j\).
This is exactly the nonsplitting required by
\texttt{IsThickMod.nonsplit} and completes the proof.
\end{proof}

\paragraph{Scope of the remaining arithmetic transfer.}
The result above constructs the actual parity presentation
and Sharp ideal.  It does not identify an endpoint block
with an old rational two-dimensional automorphic point.
At an even prime, twisting the endpoint block alone by
the fixed Teichmuller lift of \(\chi^{-j}\) removes its
finite residual cubic factor; this is a statement about
that block, not a proof that the deformed middle
character is one or that every old local and global
two-dimensional hypothesis is satisfied.  Twisting
the whole three-dimensional representation instead
changes the middle character.  Those distinct maps
must not be interchanged in a propagation argument.

Whenever a subsequent step removes finite-order local
ratios, its excluded locus is the all-pair normalized
collision locus of AC-G101.  The empty ordinary
adjacent equality locus is not adequate for that step.
The present theorem does not supply a native arithmetic
N4 assertion, a prescribed support curve, or the
remaining arithmetic adjacent-sector dimension bound.
It does provide the component chains and the exact
algebraic dichotomy to which a correctly scoped point
and propagation theorem can be applied.

\paragraph{Source and categorical dictionary.}
The exact pinned source interfaces are
\texttt{Integral.Space.Loc.EPres} and
\texttt{Integral.Space.Loc.Sharp}; the applied proofs
are \texttt{SumLe.epres}, \texttt{EpGflag.ep\_gflag},
\texttt{EpGlift.ep\_glift}, \texttt{EpMazur.ep\_mazur},
\texttt{Oddm.oddm}, \texttt{Detsec.detsec},
\texttt{ThickOfDer.thick}, and
\texttt{ChainedOfComap.sharp3}, at commit
\begin{center}
\texttt{7cbf7a0c32fca1badf05931879a37c3f8a2b03b4}.
\end{center}
Their complete containing proofs were read; exact URLs,
bytes and hashes are retained in the companion manifest.
The source's one-character dual hypothesis and its
\texttt{gpos} input were not silently transferred.
The former is replaced by the actual calculation in
AC-G101 and the latter by the explicit \(g=0\) case above.
The generic incidence radical and formal completion
arguments are spelled out from \texttt{GDetrad},
\texttt{GHomprime} and \texttt{GDetps}; the connectedness
input is the directly read Stacks result cited above.
No Lean verification or upstream axiom audit is claimed.

In scheme language, the sign action is the action of
the order-two diagonal subgroup of the grading torus;
its fixed quotient is cut out by the odd entries.
The derivation (G105.15) is a nonzero relative normal
tangent to that fixed quotient at a determinant
component.  Equation (G105.16) expresses that tangent
as an extension of the two invariant summands.
These are statements about the actual formal
representation scheme, its ideals and coherent
modules.  They require no assertion that a conjectural
motivic realization of the representation scheme
has already been constructed.
